% Technical bootstrap note for the pure newest-or-zero raw guarded calculus.
\section{Deleting the literal naming exception on raw guarded terms}
\label{sec:raw-guarded-bootstrap}

Let $\mathcal G_\circ$ be the incoming deterministic guarded calculus
with its naming-pattern exception deleted.  Its introduction witness is
always the newest ambient numerical parameter, or $0$ in an empty scope.
Its terms are raw $0,S,I$ terms; its canonical atomic paired matrices are
$I(y,0;u.h)=I(y,0;v.k)$; its other rules are the two computation schemas,
Use, full formula Transport, and Open.  The introduction mode is progressive
at an equation and by cases at a compound matrix.  Numerical scopes are
kept distinct from lists of retained clauses.

The difficulty is not soundness.  It is bootstrapping names without silently
using renaming or cut for finished $\mathcal G_\circ$ derivations.  The
following supplies a noncircular route.

\subsection{Hereditary proof families}

Write $\Gamma\Vdash A$ when, for every injective renaming of its displayed
parameters into an arbitrary larger ordered scope and every clause context
containing the renamed $\Gamma$, there is an effectively regenerated
$\mathcal G_\circ$ proof of the renamed $A$.  This is a metatheoretic
property of a proof family, not an extra rule in the checker.

Hereditary proofs can be cut into arbitrary finished target proofs.  At each
Use of the ordinary assumption being removed, regenerate the hereditary
donor in the exact numerical scope and clause context of that leaf.
Every other target proof node is left in precisely the same numerical scope,
so every selected newest parameter remains unchanged.  Induction on the
finite receiving proof proves this replacement theorem.  This argument
does not assume that the receiving proof itself can be renamed.

The algebraic rules preserve hereditary families.  An introduction whose
witness coordinate is absent from its matrix also does: the value chosen
for that coordinate has no effect.  Open is handled by allocating a fresh
parameter in each regenerated proof.  These observations permit the
padding and equation-implication constructions below to be carried out in
hereditary families.

\subsection{The arithmetic bootstrap before names}

Put $G_y(s=t)=(C(y,0,s)=C(y,0,t))$ and
$L(E)=D_{a,y}G_yE$, with $a,y$ fresh for $E$.
The coordinate $a$ is unused.  Uniformly and hereditarily,
\[
 L(E)\dashv\vdash E.
\]
Forward introduction uses the available $E$ to prove the successor row
$G_{Sz}E$ and ignores the local progress hypothesis.  Reverse elimination
opens $L(E)$, supplies its computationally true zero row, and uses its
clause at $1$.  These constructions never need a chosen numerical witness.

For equations $U=(s=t)$ and $V=(r=q)$ define
\[
 U\Rightarrow_g V
  :=D_{a,w}L\bigl(C(w,s,r)=C(w,t,q)\bigr).
\]
Again $a$ is unused.  This compound matrix admits hereditary implication
introduction and elimination by its zero and successor endpoints.
Likewise, hereditary zero/successor cases are available for a compound
formula by an unused-coordinate case introduction, Open, and Use; for an
equation apply that construction to its padding $L(E)$.

There is one important qualification.  Do not infer a generic substitution
theorem for target derivations from these constructions.  Instead let
$\mathcal E$ be ordinary equational reasoning with the two iterator
equations and literal-equation induction.  Prove the following by induction
on a \emph{source} $\mathcal E$ proof:
\[
 \mathcal E:\Gamma\vdash E
 \quad\Longrightarrow\quad
 \Gamma\sigma\Vdash E\sigma
 \quad\text{for every numerical term substitution }\sigma.
 \tag{E-embedding}
\]
Source contexts here contain equations.  A source equation is represented
as the usual constant clause in the target.

The computation and equality cases follow from hereditary algebraic
reasoning and hereditary cut.  At a source equation-induction step, write
the target-substituted invariant as $E(y)$ and put $J(y)=G_yE(y)$.
The proper source step subproof can be regenerated by the induction
hypothesis at $y=0$ and at $y=Sw$; no finished target proof is substituted.
The first gives $\Gamma\sigma,E(0)\Vdash E(1)$ and the second gives
$\Gamma\sigma,E(Sw)\Vdash E(SSw)$.  All fixed outer parameters remain
parameters of these displayed formulas; only the source induction variable
is assigned the indicated new term.  Every local entailment in this
construction retains that same $\Gamma\sigma$.

Consequently the compound formula
\[
 \Phi(y):=J(y)\Rightarrow_g E(Sy)
\]
has hereditary proofs at $0$ and at $Sw$.  At $0$ use the source base
and regenerated source step and ignore $J(0)$.  At $Sw$, remove the
successor guard from $J(Sw)$ and use the regenerated source step.
Hereditary numerical cases yield $\Phi(z)$.  Its implication elimination
gives $J(z)\Vdash J(Sz)$.  Guarded progress with an unused witness,
then Open and Use, yields $E(Sw)$; numerical cases with the original
base yield the desired $E(t)$.

All recursive calls in this argument concern proper subproofs of the
source $\mathcal E$ derivation.  All target proof replacements use
already hereditary donors.  Thus (E-embedding) is well founded and
precedes any naming or general target structural theorem.

The purely equational predecessor construction of
Section~\ref{sec:equational-predecessor} now supplies hereditary proofs,
under every numerical substitution, of
\[
 P(0)=0,\qquad P(St)=t.
\]
In particular, applying $P$ and its two instantiated equations converts
$Sb=St$ to $b=t$ in every extended target scope.

\subsection{A name generated from the newest parameter}

Define the canonical guarded packet
\[
 A(r):=D_{x,v}\bigl(
 C(v,0,C(r,0,Sx))=C(v,0,r)\bigr).
\]
The following two observations construct a hereditary proof of the
\emph{whole} outer closure $D_{a,r}A(r)$, where $a$ is unused.

First, $A(0)$ can be introduced at whichever witness the current scope
selects: its successor row computes to $0=0$, independently of the
witness.  Second, in the branch of the outer compound introduction, its
step parameter $z$ has just been appended and is newest.  Introduce
$A(Sz)$ at that parameter.  Its atomic introduction is progressive;
its successor row computes to $Sz=Sz$, independently of the progress
hypothesis.  The outer introduction therefore succeeds while ignoring
its local assumption $A(0)$.

This argument must be applied to the whole closure construction.  A
standalone proof of $A(Sz)$ is not claimed hereditary: moving $z$ to an
older position can change the selected witness.  Regenerating the whole
outer proof always allocates $z$ freshly at the required position.

Open the outer closure, supply its uniformly provable base $A(0)$, and
use its clause at $St$.  This yields $A(St)$ in any ambient scope.
Open that packet at a fresh $b$.  Its zero row is a computation theorem;
Use at $1$ and the successor computations give $Sb=St$.  The
hereditary predecessor equations give $b=t$.  This extraction can be
regenerated below any later step parameter: the name clause stays in
the context and all other ingredients are hereditary.

We may also recover the incoming paper's literal naming theorem
\[
 D_{x,y}N_t(x,y),\qquad N_t(x,y)=C(y,0,x)=C(y,0,t).
\]
Open $A(St)$ at newest $b$, extract $b=t$, and introduce $D N_t$.
Its selected witness is now exactly $b$, and the successor target follows
algebraically from $b=t$.  Close the opening.  The entire recipe is
uniform under every renaming and scope extension, because it regenerates
its own fresh witnesses; it uses no literal pattern recognizer.

\subsection{Completeness after the bootstrap}

At this point the structural proof of the incoming deterministic-introduction
paper applies with no naming-pattern case.  Induct on an original target
derivation, universally over parameter/context embeddings.  When the
default selected witness changes, uniformly generate a name for the image
of the old witness, open it at newest $b$, and apply the induction
hypothesis to the original node's proper branch in the enlarged scope.
Algebraic Transport converts between the old and new branch assumptions;
restricted cut removes the extra old assumption.  Introduce at newest $b$
and close its name.  All recursive calls concern proper original
subderivations, while the new naming proofs are generated independently.

General ordinary-assumption cut then follows by the usual fixed-proof leaf
replacement using this established structural theorem.  Finally translate
an unrestricted guarded introduction at $t$ by naming $t$ with a fresh
$b$, transporting its branch to $b$, applying the deterministic
introduction at newest $b$, and discharging the name.  Mode selection is
unchanged because it depends only on the matrix root.

Hence $\mathcal G_\circ$ and the unrestricted raw guarded calculus
derive exactly the same canonical sequents.  Composing the existing
guarded/native/iPRA interpretations gives the same full-domain
bi-interpretation.  The concrete axiomatization change is deletion of the
literal-name recognizer and its exceptional witness selection; neither
normalization nor a replacement arithmetic primitive is introduced.
